### Paragraph 1 — Main Findings: Swarm Algorithms in Precision Agriculture

The comparative analysis of the $n_{primary} = 26$ primary studies in the corpus (out of $n=33$ total included studies) reveals that swarm intelligence (SI) algorithms —specifically Particle Swarm Optimisation (PSO) and Ant Colony Optimisation (ACO)— are the dominant metaheuristic paradigms for multi-UAV path planning in precision agriculture, accounting for 30.3\% (10 of 33) and 9.1\% (3 of 33) of the literature respectively. Studies targeting monitoring and spraying applications report that PSO variants achieve average energy reductions of 15--34\% compared to systematic sweep (lawnmower) strategies in static open-field environments, where the algorithm's cognitive-social velocity dynamics $\mathbf{v}_i^{(t+1)} = \omega \mathbf{v}_i^{(t)} + c_1 r_1 (\mathbf{p}_{best,i} - \mathbf{x}_i^{(t)}) + c_2 r_2 (\mathbf{g}_{best} - \mathbf{x}_i^{(t)})$ enable global exploration of the field with minimal coverage redundancy \cite{chen2021,lin2022,liu2021cassa}. In contrast, for high-density spraying tasks or monitoring in orchards with fixed and emergent obstacles, ACO-based approaches demonstrate a 8--21\% superior route efficiency over PSO, as the virtual pheromone mechanism $\tau_{ij}$ provides a more robust framework for discrete waypoint optimisation and adaptive obstacle avoidance in complex 3D agricultural environments \cite{mathew2021,phung2021,pan2022}.

### Paragraph 2 — Contextualisation within the State of the Art and Contributions

The synthesis of the 2021--2024 corpus demonstrates a significant shift from generic path planning towards agricultural-specific optimisation models. While literature prior to 2018 often neglected the energy dynamics of small-scale multirotors, 50.0\% of the current primary studies (13 of 26) incorporate quantitative energy metrics, with a growing subset integrating the **Mulgaonkar energy model** ($P_{hover} = \sqrt{(mg)^3 / (2\rho \pi R^2)}$) to refine autonomy estimates for real-world agricultural payloads \cite{shafiq2022,xu2022,yu2021}. Furthermore, the identification of a **communication overload index** (COI) in seven studies \cite{ahmed2021,deng2023,hu2025} reveals that inter-agent coordination for swarms larger than 15 UAVs can account for 12--23\% of the mission's energy budget, a constraint that becomes the dominant bottleneck in commercial-scale field mapping (100--1,000\,ha) requiring high-frequency positional synchronisation.

### Paragraph 3 — Methodological Limitations and Future Research

Despite these algorithmic advances, the corpus reveals a persistent and critical **hardware-in-the-loop validation gap** (HIL gap): 100\% of the 26 primary studies rely exclusively on simulation, with no evidence of real-world hardware validation in agricultural field conditions. This reliance on purely software-based models leads to an unquantified "optimism bias," where simulation energy savings (mean 28.4\%) significantly exceed the conservative estimates (17.6\%) reported in the few studies using high-fidelity hardware-coupled environments (e.g., Gazebo-PX4) \cite{li2023,rao2024}. A second structural limitation is the **environmental simplification gap**: 75.8\% of studies (25 of 33) do not model dynamic obstacles, and 60.6\% (20 of 33) neglect wind and weather perturbations, which are the primary sources of mission failure in real-world spraying and monitoring. Resolving these gaps requires the adoption of the proposed **AgriSwarm-Bench** framework, emphasizing standardised energy metrics, experimentally validated aerodynamic models, and multi-agent scalability testing to bridge the transition from algorithmic theory to practical agricultural deployment.
